%\newsection{Kapitel1}

%Hier den Inhalt mit \input einfügen und die folgenden Zeilen entfernen


%%%%%%%%%%%%%%%%%%%%%%%%%%%%%%%%%%%%%%%%%%%%%%%%% Anfang - Induktivität %%%%%%%%%%%%%%%%%%%%%%%%%%%%%%%%%%%%%%%%%%%%%%

\section{Induktivität und Spule}

\s{
	Das Gegenstück zum Kondensator ist die Spule, oder genauer gesagt: Das Gegenstück zur Kapazität $C$ ist die Induktivität $L$.
	Die Induktivität ist ebenfalls ein Energiespeicher, wobei der Unterschied in der Art der genutzten Felder liegt.
	Während bei der Kapazität die Energie in Form eines elektrischen Feldes gespeichert wird,
	geschieht dies bei der Induktivität über das magnetische Feld.}
%%%%%%%%%%%%%%%%%%%%%%%%%%%%%%%%%%%%%%%%%%%%%%%%% Anfang - Lernziele Induktivität %%%%%%%%%%%%%%%%%%%%%%%%%%%%%%%%%%%%%%%%%%%%%%
	% \begin{Lernziele}{Induktivität und Spule}
	% 	Die Studierenden können
	% 	\begin{itemize}
	% 		\item zwischen Induktivität und Spule differenzieren.
	% 		\item die wichtigsten Parameter rund um die Induktivität und die Spule benennen.
	% 		\item die Induktivität einer Spule berechnen.
	% 	\end{itemize}
	% \end{Lernziele}
% \b{
% \begin{frame}
% 	\ftx{Lernziele}
% 	\begin{Lernziele}{Induktivität und Spule}
% 		Die Studierenden können
% 		\begin{itemize}
% 			\item zwischen Induktivität und Spule differenzieren.
% 			\item die wichtigsten Parameter rund um die Induktivität und die Spule benennen.
% 			\item die Induktivität einer Spule berechnen.
% 		\end{itemize}
% 	\end{Lernziele}
% \end{frame}
% }
\v{
	\begin{frame}
		\ftx{Lernziele}
		\begin{Lernziele}{Induktivität und Spule}
			Du kannst
			\begin{itemize}
				\item zwischen Induktivität und Spule differenzieren.
				\item die wichtigsten Parameter rund um die Induktivität und die Spule benennen.
				\item die Induktivität einer Spule berechnen.
				\item die Lade- und Entladeeigenschaften erläutern.
			\end{itemize}
		\end{Lernziele}
		\speech{04_Induktitaet_und_Spule1}{1}
		{In diesem Abschließenden Kapitel behandeln wir die Induktivität als Eigenschaft und die Spule als reales Bauteil.
			Am Ende dieses Kapitels kannst du zwischen Induktivität und Spule unterscheiden, kennst die wichtigsten Parameter einer Spule und der Induktivität sowie ihr Verhalten, kannst die Induktivität einer Spule berechnen und abschließend die Lade- und Entladeeigenschaften erläutern.}
	\end{frame}
}
%%%%%%%%%%%%%%%%%%%%%%%%%%%%%%%%%%%%%%%%%%%%%%%%% Ende - Lernziele Induktivität %%%%%%%%%%%%%%%%%%%%%%%%%%%%%%%%%%%%%%%%%%%%%%

\subsection{Die Induktivität $L$} \label{sec:Induktivität}
\s{
	Die Induktivität ist die Fähigkeit eines Objektes, elektrische Energie in Form eines magnetischen Feldes zu speichern.
	Diese Eigenschaft kommt zustande, sobald ein elektrischer Strom durch einen Leiter fließt um welchen sich entsprechend ein magnetisches Feld ausbildet.
	Das Formelzeichen der Induktivität ist $L$ (benannt nach dem französischen Physiker Heinrich Lenz), während die Einheit in Henry ($\mathrm{H}$) angegeben wird.

	Das einfachste und geläufigste Beispiel, an dem die Induktivität erklärt werden kann, ist die Spule.
	Obwohl sich das Phänomen der Induktivität bereits bei einem einfachen Leiter ausbildet, wird dieser Effekt durch den Aufbau einer Spule bewusst verstärkt.
	In Abbildung \ref{fig:Spule} ist ein Zylinderluftspule mit ihrer magnetischen Flussdichte $\vec{B}$ dargestellt.
}
%%%%%%%%%%%%%%%%%%%%%%%%%%%%%%%%%%%%%%%%%%%%%%%%% Anfang - Grafik Spule ohne Kern %%%%%%%%%%%%%%%%%%%%%%%%%%%%%%%%%%%%%%%%%%%%%%
% Spule ohne Kern
\begin{frame}
	\ftx{Magnetisches Feld einer Zylinderluftspule}
	\begin{figure}[H]
		\centering
		\includesvg[width=0.6\textwidth]{Magnetisches_Feld_Spule}
		\s{\caption{\textbf{Darstellung des magnetischen Feldes einer Zylinderluftspule.} Es illustriert die Feldlinien sowohl innerhalb als auch außerhalb der Spule und veranschaulicht, wie sich das Magnetfeld entlang  der Achse der Spule konzentriert und außerhalb der Spule abnimmt.}}
		\alttext{Darstellung des magnetischen Feldes einer Zylinderluftspule. Links ist eine dreidimensionale Spule gezeigt; ein Pfeil markiert den Stromfluss I durch die Wicklung. Rechts ist das zugehörige Magnetfeld als Feldlinienbild dargestellt: Im Inneren verlaufen die Feldlinien dicht und parallel zur Spulenachse, außerhalb schließen sie sich in großen Bögen um die Spule. Symbole mit Punkten und Kreuzen kennzeichnen Stromrichtung in und aus der Bildebene.}
		\label{fig:Spule}
	\end{figure}
	\speech{04_Induktitaet_und_Spule2}{1}
	{Ähnlich wie eine Kapazität hat eine Induktivität die Fähigkeit Energie im Feld zu Speichern. Sie speichert aber nicht wie die vorher kennengelernte Kapazität ihre Energie im Elektrischen Feld sondern im magnetischen Feld, was durch den Stromfluss entsteht. Jeder Strom der fließt erstellt ein Magnetfeld zirkular um den Leiter herum, wie in dem rechten Bild zu sehen ist. Jeder Leiter in dem Querschnitt der Spule erstellt ein zirkulares Magnetfeld, welches eine bestimmte orientierung hat, abhängig von der richtung des Stromes.}
\end{frame}
%%%%%%%%%%%%%%%%%%%%%%%%%%%%%%%%%%%%%%%%%%%%%%%%% Ende - Grafik Spule ohne Kern %%%%%%%%%%%%%%%%%%%%%%%%%%%%%%%%%%%%%%%%%%%%%%
\s{
	Ihren Namen verdankt diese Spule ihrer zylindrischen Form in Längsrichtung und des Mediums Luft im inneren ihrer Wicklungen.
	Wenn nun ein Strom durch die Spule fließt, überlagern sich die Magnetfelder der einzelnen Wicklungen zu einem stärkeren Magnetfeld. Die Stärke hängt dabei von der Windungszahl
	und der Stromstärke, sowie der Entfernung zum Leiter, ab. Je nach Orientierung des Stromes und der Wicklungen, entsteht bei einer Zylinderluftspule auf der einen Seite ein Nord- und auf der anderen Seite ein Südpol. 
	
	Das Thema der magnetischen Größen wird ausführlich in Modul 6 behandelt. Zum Verständnis der Funktionsweise einer Spule und der daraus resultierenden Induktivität, soll eine kurze Einführung zur Entstehung von Magnetismus helfen. Sobald Strom durch einen geraden Leiter fließt, entsteht ein magnetisches Feld konzentrisch um den Leiter herum. Abhängig von Stromstärke und Stromrichtung ändert sich auch das magnetische Feld. Die Rechte-Faust-Regel dient dabei als Gedankenstütze zur Bestimmung der Richtung und Orientierung des Magnetfeldes.
}
%%%%%%%%%%%%%%%%%%%%%%%%%%%%%%%%%%%%%%%%%%%%%%%%% Anfang - Grafik Lenzsche Regel %%%%%%%%%%%%%%%%%%%%%%%%%%%%%%%%%%%%%%%%%%%%%%
\begin{frame}
	\ftx{Die Rechte-Hand-Regel}
	\begin{figure}[H]
		\centering
		\includesvg[width=1\textwidth]{Rechte-Faust-Regel}
		\s{\caption{\textbf{Rechte-Hand-Regel zur Bestimmung der Richtung des Magnetfeldes.} Ein stromdurchflossenen Leiter erzeugt ein Magnetfeld das mit
				hilfe der Rechten-Hand-Regel bestimmt wird. Wenn der Daumen der rechten Hand in die Richtung des Stroms zeigt, dann krümmen sich die übrigen Finger in
				die Richtung der Feldlinien des Magnetfeldes.}}
			\alttext{Illustration der Rechte-Hand-Regel an einem stromdurchflossenen geraden Leiter. Ein schräg verlaufender Leiter ist mit einem Pfeil für den Strom I nach oben dargestellt. Um den Leiter ist eine ringförmige Feldlinie des Magnetfeldes B als Pfeil eingezeichnet. Daneben zeigt eine rechte Hand: Der Daumen weist in Stromrichtung, die gekrümmten Finger geben die Umlaufrichtung der Magnetfeldlinien an.}
		\label{fig:Rechtehand}
	\end{figure}
	\speech{04_Induktitaet_und_Spule3}{1}
	{Die Richtung des Magnetfeld kann anhand der rechten hand regel Festgestellt werden. Der Daumen zeigt in Richtung des Stromes, und die krümmung der Finger zeigt den Verlauf des Magnetfeldes.
		Wichtig ist, die Felder e beeinflussen sich nicht gegenseitig und lassen sich überlagern. Das heißt eine Addition, aber auch Auslöschung ist möglich.
		Wir haben nun gelernt, das ein Strom ein magnetisches Feld erzeugt. Manche kennen es aber auch, wir können durch ein Zeitlich veränderliches Magnetfeld eine Spannung in Leitern induzieren.
	}
\end{frame}
%%%%%%%%%%%%%%%%%%%%%%%%%%%%%%%%%%%%%%%%%%%%%%%%% Ende - Grafik Lenzsche Regel %%%%%%%%%%%%%%%%%%%%%%%%%%%%%%%%%%%%%%%%%%%%%%

\s{
	Entsprechend der Abbildung \ref{fig:Rechtehand} verläuft das Magnetfeld in Richtung der Fingerspitzen konzentrisch um den Leiter herum, sofern man den Daumen in Richtung der technischen Stromrichtung,
	also entgegen der Fließrichtung der Elektronen, ausrichtet. Das Magnetfeld dreht sich von Süd nach Nord und hat eine Flussdichte $\vec{B}$, sowie eine Feldstärke $\vec{H}$.
	Dabei ist zu beachten, dass sich die magnetische Flussdichte $\vec{B}$ analog zur elektrischen Feldstärke $\vec{E}$ , und die magnetische Feldstärke $\vec{H}$ analog zur elektrischen Flussdiche $\vec{D}$, verhält.
	Im Modul 6 wird das Thema Magnetismus ausführlicher behandelt. In diesem Modul geht es primär um das Kennenlernen der Spule und der Induktivität.
}

%%%%%%%%%%%%%%%%%%%%%%%%%%%%%%%%%%%
\begin{frame}[t]
	\ftx{Die Induzierte Spannung}
	\b{
			\begin{itemize}\centering
				\item $ \displaystyle U_\mathrm{i} \sim \frac{\mathrm{d}\phi}{\mathrm{d}t}$\vspace{0.5cm}
			
		\uncover<2->{\centering
			
				\item  $ \displaystyle U_\mathrm{i}= -N \cdot \frac{\mathrm{d}\phi}{\mathrm{d}t}$\vspace{0.5cm}
		
		}
		\uncover<3->{\centering
			
				\item $ \displaystyle U_\mathrm{i} \sim \frac{\mathrm{d}I}{\mathrm{d}t}$\vspace{0.5cm}
		
		}
		\uncover<4->{\centering
		
				\item $\displaystyle U_\mathrm{i}= -L \cdot \frac{\mathrm{d}I}{\mathrm{d}t}$\vspace{0.5cm}
		}
	\end{itemize}
	}
	\speech{04_Induktitaet_und_Spule4}{1}
	{Die Spannung U-i die in einer vom magnetfeld durchsetzten Spule induziert wird ist proportional abhängig von der zeitlichen ableitung des magnetischen Flusses. Das bedeutet, je schneller sich das magnetfeld verändert, umso höher ist die Spannung, die in der Leiterschleife induziert wird.}
	\speech{04_Induktitaet_und_Spule4}{2}
	{Wenn wir nun diese Leiterschliefe mehrfach haben, also N aneinandergschlossene Leiterschleifen, wie bei einer gewickelten Spule, dann wird in jeder einzelnen Spule diese Spannung induziert in die Spannung erhöht sich um den Faktor N, der Anzahl der Wicklungen.
	}
	\speech{04_Induktitaet_und_Spule4}{3}
	{Die Ursache Des magnetfeld ist aber immer noch der Strom durch den Leiter, was bedeutet, das die induzierte Spannung in dem Leiter auch proportional zu dem Strom durch den leiter ist.
	}
	\speech{04_Induktitaet_und_Spule4}{4}
	{Dieser Proportionalitätsfaktor zwischen dem Strom und der induzierten Spannung nennt sich Induktivität und wird mit dem Buchstaben L abgekürzt.
	Mit anderen Worten, jeh höher die die frequenz des Stromes ist umso größer ist die Spannung die an der induktivität abfällt beziehungweise induziert wird.
	}
\end{frame}

\s{
	Der britische Physiker und Chemiker Michael Faraday führte Mitte des 19. Jahrhunderts Experimente durch,
	bei denen er eine Spule in der Nähe eines Magneten bewegte oder den Strom in einer benachbarten Spule veränderte.
	Dabei stellte er fest, dass sich die Größe der induzierten Spannung proportional zur Änderungsrate des Magnetfeldes verhält.

	\begin{equation}
		U_\mathrm{i} \sim \frac{\mathrm{d}\phi}{\mathrm{d}t}
	\end{equation}

	Faraday stellte fest, dass der Proportionalitätsfaktor von der Anzahl der Wicklungen der Spule abhängt.

	\begin{equation}
		U_\mathrm{i}= -N \cdot \frac{\mathrm{d}\phi}{\mathrm{d}t}
	\end{equation}

	Etwa zur gleichen Zeit entdeckte der amerikanische Physiker Joseph Henry das Phänomen der Selbstinduktion,
	wobei ein elektrischer Strom in einer Spule eine Spannung in der Spule selbst induzieren kann, wenn der Strom verändert wird.
	Auch hier wurde ein proportionales Verhalten festgestellt.

	\begin{equation}
		U_\mathrm{i} \sim \frac{\mathrm{d}I}{\mathrm{d}t}
	\end{equation}

	Diese Eigenschaft der Spule die durch das Speichern magnetischer Energie entsteht ist die Induktivität $L$,
	deren Einheit dem Physiker zu Ehren als Henry $[\mathrm{H}]$ benannt wurde.

	\begin{equation}
		U_\mathrm{i}= -L \cdot \frac{\mathrm{d}I}{\mathrm{d}t}
	\end{equation}

	Das Minus-Zeichen kommt daher, dass der induzierte Strom nach der Lenz'schen Regel seiner Ursache entgegenwirkt.
}% nur im Skript


%%%%%%%%%%%%%%%%%%%%%%%%%%%%%%%%%%%%%%%%%%%%%%%%% Anfang - Formeln Zeitabhängigkeit Induktivität	 %%%%%%%%%%%%%%%%%%%%%%%%%%%%%%%%%%%%%%%%%%%%%%
%%%%% Zeitabhängige Formeln für induktivität

\begin{frame}[t]
	\begin{overlayarea}{\textwidth}{0.9\textheight}
		\ftx{Zeitabhängige Formeln der Induktivität}
		\b{
			
			\begin{itemize}
				\item Spannung über der Induktivität
			\end{itemize}
			\begin{equation*}
				u_\mathrm{L}(t) = L \cdot \frac{di_\mathrm{L}(t)}{dt}
			\end{equation*}
			\only<2->{
				\begin{itemize}
					\item Strom durch die Induktivität
				\end{itemize}
				\begin{equation*}
					i_\mathrm{L}(t) = \frac{1}{L} \cdot \int u_\mathrm{L}(t) \, dt
				\end{equation*}
			}
			\only<3->{
				\vspace{-1cm}
				\begin{Merksatz}{Spule}

					\begin{itemize}
						\item Für DC stellt die Induktivität einen Kurzschluss dar
						\item Der Strom durch eine Induktivität kann sich nicht sprunghaft ändern
						\item Bei plötzlichen Stromänderungen (Schalten) entstehen sehr hohe Spannungen
					\end{itemize}
				\end{Merksatz}

			}
		}
	\end{overlayarea}
	\speech{04_Induktitaet_und_Spule5}{1}
	{Hier nochmal zusammengefasst, die Spannung über der Induktivität ist abhängig von dem Induktivitätswert L und der zeitlichen Änderung des Stromes.}
	\speech{04_Induktitaet_und_Spule5}{2}
	{Diese obere Formel können wir umformen, indem wir sie zeitlich integrieren und durch den Faktor L teilen, und somit den Strom durch die induktivität in abhängigkeit von einer anliegenden Spannung erhalten.}
	\speech{04_Induktitaet_und_Spule5}{3}
	{Es ist an diesen zwei formeln einfach zu sehen, das an einer induktivität nur eine Spannung abfällt, wenn wir einen zeitlich veränderlichen Strom haben. Wenn du also nur DC Strom durch eine Induktivität schickst, dann werden Null Volt über die induktivität abfallen.
		Der Strom durch eine induktivität kann sich nicht sprunghaft ändern, da das in einer unendlichen Spannung über der induktivität resultieren würde. Als vergleich betrachten wir das integral über die Spannung über der induktivität. Da zeigt sich das gleiche Bild. Denn für eine Sprunghafte änderung des Stromes durch eine induktivität müsste das zeitliche integral über die Spannung ebenfalls unendlich werden.
		Das zeigt aber auch, das beim schalten von Strömen an induktivitäten sehr hohe spannungen entstehen können, die die Schaltung in der sie sich befinden beschädigen oder zerstören können.
	}
\end{frame}

\s{
	Analog zum Kondensator, kann auch das zeitabhängige Verhalten für die Spule folgendermaßen ausgedrückt werden. Die Spannung über der Induktivität ist
	\begin{equation}
		u_\mathrm{L}(t) = L \cdot \frac{di_\mathrm{L}(t)}{dt}
	\end{equation}
	und der Strom durch die Induktivität ist
	\begin{equation}
		i_\mathrm{L}(t) = \frac{1}{L} \cdot \int u_\mathrm{L}(t)\,dt
	\end{equation}
}

%%%%%%%%%%%%%%%%%%%%%%%%%%%%%%%%%%%%%%%%%%%%%%%%% Ende - Formeln Zeitabhängigkeit Induktivität	 %%%%%%%%%%%%%%%%%%%%%%%%%%%%%%%%%%%%%%%%%%%%%%

\subsection{Die Spule als Bauelement}

\s{
	Mit dem Wissen über das magnetische Feld kann nun der Aufbau und die Funktionsweise der Spule erklärt werden. Ein einfaches Beispiel dafür ist die Zylinderluftspule.
	Wird ein Leiter zu einer Leiterschleife geformt, bildet sich das Magnetfeld so aus, dass es durch die Schleifenöffnung hindurchgeht.
	Werden nun mehrere Leiterschleifen hintereinander gewickelt, wird das Magnetfeld um den Faktor der Windungszahl verstärkt.
	Diese Anordnung von gewickelten Leiterschleifen anhand von realen Bauteilen ist in Abbildung \ref{fig:Vielespulen} dargestellt.
}

\begin{frame}[t]
	\ftx{Die Spule als Bauelement}
	\b{
		\begin{itemize}
			\item Verschiedene Arten und Bauformen von Spulen
		\end{itemize}}
	\begin{figure}[H]
		\centering
		\includegraphics[width=0.7\textwidth]{sp}
		\s{\caption{\textbf{Verschiedene Arten und Bauformen von Spulen.} Darunter Zylinderspulen, Ringkernspulen, Eisenkernspulen und Luftspulen. Eine größere Anzahl von Windungen erhöht die Konzentration des Magnetfelds und steigert die Induktivität der Spule. Ein Eisenkern kann zusätzlich verwendet werden, um die Magnetfeldlinien zu bündeln und die Effizienz zu verbessern.}}
\alttext{Die Abbildung zeigt ein Foto verschiedener Spulen und Induktivitäten in unterschiedlichen Bauformen auf einer hellen Oberfläche. Zu sehen sind unter anderem eine Ringkernspule, flache und zylindrische Luftspulen und bedrahtete Induktivitäten sowie eine große flache Spule. Die Bauteile sind in mehreren Reihen angeordnet und zeigen den Größen- und Bauformvergleich von Spulen.}
		\label{fig:Vielespulen}
	\end{figure}
	\speech{04_Induktitaet_und_Spule6}{1}
	{Um eine Induktivität nachzuahmen werden Spulen verwendet. Diese sind in ihrer grundform zylindrisch gewickelte Leiter, welche das magnetfeld vor allem im inneren haben. Um die induktivität zu erhöhen werden permeable materialien verwendet, wie oben links in gelb zu sehen ist. Diese Form der Spule wird häufig als filter benutzt, um hochfrequente störungen zu blocken. Es gibt aber auch andere Spulen, mit denen energie Übertragen werden kann wie mit den Flachen Spulen in der Mitte die für NFC oder ähnliche kommunikationen verwendet werden können.
	Da eine induktivität an sich keine verluste besitzt, da sie die energie in ihrem Feld speichert und abgibt, sind die Verluste primär durch den ohmschen Widerstand der Spule und magnetische Verluste durch den verwendeten Kern der Spule ergeben, sind Spulen für große Ströme auch räumlich größer. Die Spule mit dem gelben Kern wird deutlich höhere Ströme abkönnen als die kleine t-h-t Spule auf der rechten seite. Ebenso wird die untere runde Spule für kontaktloses schnelladen verwendet, und die obere nur für kleinere Leistungen wie kontaktloses Bezahlen verwendet.}
\end{frame}
\s{
	Was alle Spulen verbindet, ist ihre Fähigkeit, die Eigenschaft der Induktivität nutzbar zu machen,
	wie es im folgenden Schaltbilder \ref{fig:Induktivität_L_EU} und \ref{fig:Induktivität_ui} dargestellt ist. Diese Schaldbilder repräsentieren jedoch nur
	die idealisierte, nutzbare Induktivität nicht das realle Bauteil.
}% nur im Skript


%%%%%%%%%%%%%%%%%%%%%%%%%%%%%%%%%%%%%%%%%% Anfang - Schaltzeichen Induktivität EU, US	 %%%%%%%%%%%%%%%%%%%%%%%%%%%%%%%%%%%%%%%%%%%%%%
% Europäisches und amerikanisches Symbol
\begin{frame}[t]
	\ftx{Die ideale Spule}
	\b{
		\begin{itemize}
			\item Europäisches (links) und amerikanisches (rechts) Schaltelement
		\end{itemize}
	}
		\begin{figure}[H]
			\centering
			\input{Tikz/src/04_Induktivitaet_Spule/Spule_EU_AM_Schaltungselement}
			\s{\caption{\textbf{Schaltungselement für Spule.} Europäischen (links) und amerikanisches (rechts).}}
			\label{fig:Induktivität_L_EU}
		\end{figure}
	
	%%%%%%%%%%%%%%%%%%%%%%%%%%%%%%%%%%%%%%%%%%%%%%%%% Ende - Schaltzeichen Induktivität	EU, US %%%%%%%%%%%%%%%%%%%%%%%%%%%%%%%%%%%%%%%%%%%%%%
	%%%%%%%%%%%%%%%%%%%%%%%%%%%%%%%%%%%%%%%%%%%%%%%%% Anfang - Schaltzeichen Induktivität U,I	 %%%%%%%%%%%%%%%%%%%%%%%%%%%%%%%%%%%%%%%%%%%%%%
	% Ideale Spule
	\b{\uncover<2->{
		\vspace{1cm}
		\begin{itemize}
			\item Die ideale Spule als Schaltsymbol
		\end{itemize}
		\begin{figure}[H]
		\begin{center}
			\input{Tikz/src/04_Induktivitaet_Spule/Ersatzschaltbild_ideale_Spule}
		
		\end{center}
	\end{figure}}
}
\speech{04_Induktitaet_und_Spule7}{1}
	{In unseren Schaltungen markieren wir induktivitäten als ein schwarzes Rechteck. Dies stellt die ideale induktivität dar, die es in der realität leider nicht gibt.}
\speech{04_Induktitaet_und_Spule7}{2}
	{Durch die verwendeten Leiter, die ebenfalls nicht ideal sind, besitzt die induktivität einen Ohmschen Widerstand in Reihe, der die Verluste der Spule erhöhen, was vor allem bei hohen Strömen probleme bereitet. Zusätzlich stellen die gewickelten Drähte die möglichst auf kleinem raum gewickelt werden und durch eine isolationsschicht elektrisch getrennt sind eine kapazität dar, die parallel zu dem gewickelten leiter liegt. Das macht die verwendung von Spulen zum blocken von sehr hohen Frequenzen schwierig, die parallele Kapazität dominant wird, auch wenn diese häufig klein ist.}
\end{frame}

\begin{frame}
	\ftx{Die reale Spule}
		
	\b{
		\begin{itemize}
			\item Ersatzschaltbild einer realen Spule
		\end{itemize}
			\begin{figure}[H]
				\centering
				\input{Tikz/src/04_Induktivitaet_Spule/Ersatzschaltbild_reale_Spule}
			\end{figure}
		\begin{itemize}
			\item $C_\mathrm{p}$ (Parallelkapazität) und $R_\mathrm{s}$ (Serienwiderstand) beschreiben die parasitären Effekte
		\end{itemize}
	}
	
\end{frame}
%%%%%%%%%%%%%%%%%%%%%%%%%%%%%%%%%%%%%%%%%%%%%%%%% Ende - Schaltzeichen Induktivität	U,I  %%%%%%%%%%%%%%%%%%%%%%%%%%%%%%%%%%%%%%%%%%%%%%%%%%%%%

\s{
	\begin{figure}[H]
		\begin{center}
			\input{Tikz/src/04_Induktivitaet_Spule/Ersatzschaltbild_ideale_Spule}
			\caption{\textbf{Schaltzeichen einer idealen Spule.} Mit einem Spannungsabfall und einem Stromfluss.}
			\label{fig:Induktivität_ui}
		\end{center}
	\end{figure}

	Für eine realitätsgetreue Nachbildung der Schaltung ist es wichtig die Spule als Bauteil
	vollständig zu beschreiben. Dies wird mithilfe eines Ersatzschaltbildes gemacht, welches auch die
	ungewollten, also parasitären Eigenschaften des Bauteils, nachbildet. Abbildung \ref{fig:Schaltbild einer realen Spule} zeigt ein solches
	Ersatzschaltbild.

	\begin{figure}[H]
		\centering
		\input{Tikz/src/04_Induktivitaet_Spule/Ersatzschaltbild_reale_Spule}
		\caption{\textbf{Ersatzschaltbild einer realen Spule.} Die ideale Spule $L$ wird durch einen Serien-Widerstand $R_\mathrm{s}$
			und einen parallel geschalteten Kondensator $C_\mathrm{p}$ ergänzt, um parasitäre Effekte zu berücksichtigen.}
		\label{fig:Schaltbild einer realen Spule}
	\end{figure}

%%%%%%%%%%%%%%%%%%%%%%%%%%%%%%%%%%%%%%%%%%%%%%%%% Ende - ESB reale Spule	 %%%%%%%%%%%%%%%%%%%%%%%%%%%%%%%%%%%%%%%%%%%%%%
	Im Ersatzschaltbild einer Spule werden nicht nur die nutzbare Induktivität, sondern auch die
	parasitären Eigenschaften berücksichtigt, die in realen Spulen auftreten. Diese parasitären
	Eigenschaften entstehen beispielsweise durch die Kapazitäten zwischen den Wicklungen und die ohmschen Ver-
	luste im Material. Das modellierte Ersatzschaltbild ermöglicht es, diese Effekte im Schaltungsentwurf
	zu berücksichtigen und das Schaltverhalten besser vorherzusagen.
}% nur im Skript
%%%%%%%%%%%%%%%%%%%%%%%%%%%%%%%%%%%%%%%%%%%%%%%%% Anfang - Merksatz Induktivität %%%%%%%%%%%%%%%%%%%%%%%%%%%%%%%%%%%%%%%%%%%%%%
\begin{frame}
	\ftx{Merkesatz}
	\begin{Merksatz}{Spule}
	Die Spule ist der verzweifelte Versuch, eine Induktivität nachzubilden.
	\end {Merksatz}
	\speech{04_Induktitaet_und_Spule8}{1}
	{Wir merken also, die Spule ist nur der verzweifelte Versuch eine Induktivität nachzubilden.}
\end{frame}

%%%%%%%%%%%%%%%%%%%%%%%%%%%%%%%%%%%%%%%%%%%%%%%%% Ende - Merksatz Induktivität %%%%%%%%%%%%%%%%%%%%%%%%%%%%%%%%%%%%%%%%%%%%%%
%%%%%%%%%%%%%%%%%%%%%%%%%%%%%%%%%%%%%%%%%%%%%%%%% Anfang - Permeabilität	 %%%%%%%%%%%%%%%%%%%%%%%%%%%%%%%%%%%%%%%%%%%%%
\subsection{Die Permeabilität}
\s{
	Ähnlich wie im Fall des Kondensators und der Kapazität, unterscheidet sich auch die Stärke des Magnetfeldes abhängig vom Material.
	Was beim Kondensator das Dielektrikum und die Permitivität sind, wird bei der Spule als ferromagnetisches Material und Permeabilität bezeichnet.
	Ein Beispiel einer Spule mit einem ferromagnetischen Material ist die Ringkernspule in Abbildung \ref{fig:Ringkernspule}.
}
%%%%%%%%%%%%%%%%%%%%%%%%%%%%%%%%%%%%%%%%%%%%%%%%% Anfang - Grafik Ringkernspule %%%%%%%%%%%%%%%%%%%%%%%%%%%%%%%%%%%%%%%%%%%%%%
% Ringkernspule

\begin{frame}[t]
	\ftx{Die Permeabilität}
	\b{\begin{itemize}
		\item Spule mit einem ferromagnetischen Material zur erhöhunug der Permeabilität
	\end{itemize}}

	\begin{figure}[H]
		\includesvg[width=1\textwidth]{Magnetfeldline_Ringspule}
		\s{\caption{\textbf{Ringkernspule mit Eisenkern.} Der Eisenkern verstärkt das Magnetfeld und erhöht die Induktivität der Spule. Diese Bauweise minimiert Streufelder und verbessert die Effizienz, indem sie das Magnetfeld innerhalb des Kerns konzentriert.}}
		\alttext{Schematische Darstellung einer Ringkernspule mit Eisenkern. Ein ringförmiger Kern ist mit mehreren Windungen umwickelt. Unten sind zwei Anschlussdrähte dargestellt; rote Pfeile markieren die Stromrichtung I in den Leitungen. Im Inneren des Kerns zeigen orange Pfeile das Magnetfeld B, das im Ring geschlossen verläuft und im Kern konzentriert ist.}
		\label{fig:Ringkernspule}
	\end{figure}
	\speech{04_Induktitaet_und_Spule9}{1}
	{Ähnlich wie das Dielektrikum eines Kondensators, gibt es auch bei Spulen Materialien, die die induktivität bei gleichbleibender Wicklungszahl erhöhen können. Diese eigenschaft wird mit der Permeabilität des Materials beschrieben.}
\end{frame}
%%%%%%%%%%%%%%%%%%%%%%%%%%%%%%%%%%%%%%%%%%%%%%%%% Ende - Grafik Ringkernspule %%%%%%%%%%%%%%%%%%%%%%%%%%%%%%%%%%%%%%%%%%%%%%
\s{
	Bei der Ringkernspule handelt es sich um einen mit Draht umwickelten Ferritring.
	Innerhalb des Materials findet eine Ausrichtung auf atomarer Ebene statt, was zu einer Verstärkung des Magnetfeldes führt.
	Die Stärke der magnetischen Flussdichte $\vec{B}$ wird dabei neben der Dimensionen und der Stromstärke, von den Materialeigenschaften, insbesondere von der Permeabilität, bestimmt.
    Die folgende Tabelle \ref{tab:permeabilität} zeigt die Permeabilitätszahl einiger gängiger Materialien:
}
\begin{frame}[t]
	\ftx{Die Permeabilität}

\b{\begin{itemize}
	\item Relative Permeabilitätszahl verschiedener Materialien
\end{itemize}
\vspace{1cm}
}
	\begin{table}[H]
		\centering
		\resizebox{0.3\textwidth}{!}{%
			\begin{tabular}{lc}
				\toprule
				\textbf{Material}  & \textbf{\boldmath$\mu_\mathrm{r}$} \\
				\midrule
				Supraleiter 1. Art & $0$                                \\
				Vakuum             & $1$                                \\
				Luft (bei STP)     & $1.00000037$                       \\
				Kupfer             & $0.999994$                         \\
				Gold               & $0.999964$                         \\
				Aluminium          & $1.000022$                         \\
				Eisen              & $300-10000$                        \\
				Ferrit             & $4-15000$                          \\
				\bottomrule
			\end{tabular}%
		}
		\s{\caption{\textbf{Relative Permeabilitätszahl verschiedener Materialien}: Diese dimensionslose Zahl gibt an,
				wie stark das magnetische Feld in einem Material im Vergleich zum Vakuum beeinflusst wird. Sie ist ein Maß für die
				Fähigkeit des Materials, magnetische Flüsse zu speichern oder zu leiten. Materialien mit einer höheren Permeabilitätszahl
				erhöhen die Induktivität von Spulen.}}
		\alttext{Die Tabelle zeigt die relative Permeabilitätszahl müh r für verschiedene Materialien an. Aufgelistet sind: Supraleiter erster Art mit der Zahl 0, Vakuum hat ein müh r von 1, Luft bei STP hat 1,00000037, Kupfer 0,999994, Gold 0,999964, Aluminium 1,000022, Eisen 300 bis 10000 und Ferrit 4 bis 15000.}
		\label{tab:permeabilität}
		\speech{04_Induktitaet_und_Spule10}{1}
		{Die einheitenlose relative Permeabilität müh-r eines Material gibt an, wie stark das magnetische Feld im vergleich zum Vakuum beeinflusst wird. Eisen und Ferrit zum beispiel, welche häufig in Spulen verwendet werden haben sehr hohe permeabilitäten im vergleich zu Luft oder vakuum, kommen aber auch mit anderen nachteilen. Sie zeigen die Fähigkeit des Materials magnetische flüsse zu speichern oder zu leiten. }
	\end{table}
\end{frame}
%%%%% Tabelle Kästchen
% \begin{table}[H]
%     \centering
%     \begin{tabular}{|l|c|}
%         \hline
%         \textbf{Material} & \textbf{\boldmath$\mu_\mathrm{r}$} \\
%         \hline
%         Supraleiter 1. Art & $0$\\
%         \hline
%         Vakuum & $1$ \\
%         \hline
%         Luft (bei STP) & $1.00000037$ \\
%         \hline
%         Kupfer & $0.999994$ \\
%         \hline
%         Gold & $0.999964$ \\
%         \hline
%         Aluminium & $1.000022$ \\
%         \hline
%         Eisen & $300-10000$ \\
%         \hline
%         Ferrit & $4-15000$ \\
%         \hline
%     \end{tabular}
%     \caption{Permeabilitätszahl verschiedener Materialien}
%     \label{tab:permeabilität}
% \end{table}

%%%%%%%%%%%%%%%%%%%%%%%%%%%%%%%%%%%%%%%%%%%%%%%%% Ende - Permeabilität	 %%%%%%%%%%%%%%%%%%%%%%%%%%%%%%%%%%%%%%%%%%%%%%

\s{
	Zur vollständigen Beschreibung der Eigenschaften des ferromagnetischen Materials wird neben der Permeabilitätszahl
	noch die magnetische Feldkonstante $\mu_0$ benötigt. Dabei handelt es sich um das Verhalten des magnetische Feldes im Vakuum.
	Während die Permeabilitätszahl $\mu_r$ einheitslos ist, bringt die magnetische Feldkonstante die Einheit $\mathrm{H/m}$ mit.
	Sie wird mit der relativen Permeabilität des Materials multipliziert was zusammen die Permeabilitätskonstante $\mu$ ergibt.
}

%%%%%%%%%%%%%%%%%%%%%%%%%%%%%%%%%%%%%%%%%%%%%%%%% Anfang - Formeln Mü %%%%%%%%%%%%%%%%%%%%%%%%%%%%%%%%%%%%%%%%%%%%%%
\begin{frame}
	\ftx{Die Permeabilität}

	\b{
		\begin{equation*}
			\mu = \mu_\mathrm{r} \cdot \mu_\mathrm{0} \quad \quad \text{[H/m]}
		\end{equation*}

		% Erklärungen der Permeabilität
		\begin{align*}
			\mu            & : \text{Permeabilitätskonstante }\text{ [H/m]}                                 \\
			\mu_\mathrm{0} & : \text{Magnetische Feldkonstante (}\approx 4\pi \times 10^{-7} \text{ [H/m])} \\
			\mu_\mathrm{r} & : \text{Relative Permeabilität des Materials (Einheitenfrei)}                  \\
		\end{align*}
	}
	\speech{04_Induktitaet_und_Spule11}{1}
	{Sie setzt sich zusammen aus der einheitenlosen relativen Permeabilität müh-r und der magnetischen Feldkonstanten müh-null und besitzt die Einheit Henry pro meter. Diese brauchen wir um die Magnetische FLussdichte in einer Spule zu berechnen.}
\end{frame}
\s{

	\begin{equation*}
		[\mu] = 1\, \frac{\text{Henry}}{\text{m}} = 1\, \frac{\text{H}}{\text{m}} = 1\, \frac{\text{Vs}}{\text{Am}}
	\end{equation*}
	\vspace{0.1cm}
	\begin{equation}
		\mu = \mu_\mathrm{r} \cdot \mu_\mathrm{0}
	\end{equation}

	% Erklärungen der Permeabilität
	\begin{align*}
		\mu            & : \text{Permeabilitätskonstante}                                  \\
		\mu_\mathrm{0} & : \text{Magnetische Feldkonstante ($\approx 4\pi \cdot 10^{-7}$)} \\
		\mu_\mathrm{r} & : \text{Relative Permeabilitätszahl}                              \\
	\end{align*}

	%%%%%%%%%%%%%%%%%%%%%%%%%%%%%%%%%%%%%%%%%%%%%%%%% Ende - Formeln Mü %%%%%%%%%%%%%%%%%%%%%%%%%%%%%%%%%%%%%%%%%%%%%%


	Mit Hilfe der Permeabilität, des Stromes, der Windungszahl und der Länge der Spule,
	kann die magnetische Flussdichte $\vec{B}$ mit der Einheit Tesla $\mathrm{T}$ berechnet werden.
}

%%%%%%%%%%%%%%%%%%%%%%%%%%%%%%%%%%%%%%%%%%%%%%%%% Anfang - Formeln Induktivität %%%%%%%%%%%%%%%%%%%%%%%%%%%%%%%%%%%%%%%%%%%%%%
\begin{frame}[t]
	\ftx{Die magnetische Flussdichte $B$}
	\b{
		% Formel für die magnetische Flussdichte
		\begin{equation*}
			B = \frac{\mu \cdot I \cdot N} {l}
		\end{equation*}
\uncover<2->{


		% Herleitung der Einheiten
		\begin{align*}
			[B] = \frac{\text{H/m} \cdot \text{A} \cdot 1}{\text{m}}
			= \frac{\text{H} \cdot \text{A}}{\text{m}^2}
			= \frac{(\text{V} \cdot \text{s / } \colorcancel{\text{A}}{red}) \cdot \colorcancel{\text{A}}{red}}{\text{m}^2}
			= \frac{\text{V} \cdot \text{s}}{\text{m}^2}
			= \text{T (Tesla)}
		\end{align*}
}
		% Erklärung der Symbole und Herleitung der Einheiten
		\begin{align*}
			\mu & : \text{Permeabilitätskonstante} \\
			I   & : \text{Elektrische Stromstärke} \\
			N   & : \text{Anzahl der Windungen}    \\
			l   & : \text{Länge des Magnetkreises}
		\end{align*}
		\uncover<3->{

			\begin{equation*}
				\vec{B} = \mu \cdot \vec{H}
			\end{equation*}
		}
	}
	\speech{04_Induktitaet_und_Spule12}{1}
	{Die magnetische Flussdichte B, die bis jetzt betrachtet wurde für ideal leitende Leiter-spule, berechnet sich über B ist gleich müh mal I mal N geteilt durch die Länge der Spule.
	Was bedeutet, je mehr windungen wir auf die länge l bekommen und je höher der Strom durch diese  ist, umso stärker ist die magnetische Flussdichte in der Spule.}
	\silence{2}
	\speech{04_Induktitaet_und_Spule12}{2}
	{Ihre Einheit lässt sich einfach herleiten durch die permeabilitätskonstante mit Henry pro meter, mal dem Strom in Ampere, die Windungszahl N ist einheitenlos, geteilt durch die länge der Spule in Metern.
	Wir können einiges kürzen, so das am Ende Volt mal Sekunde pro quadratmeter herauskommt, was zusammengefasst wird zu der Einheit Tesla.
	Ein Tesla ist also eine Volstsekunde pro quadratmeter.
	}
	\silence{3}
	\speech{04_Induktitaet_und_Spule13}{3}
	{Während die Magnetische Flussdichte B von dem durchdringten material abhängig ist, ist die Magnetische Feldstärke H unabhängig vom Material. Erst die magnetische Flussdichte bezieht auch das verwendete Material im inneren der Spule.}
\end{frame}
\s{

	\begin{equation}
		\vec{B} = \mu \cdot \vec{H}
	\end{equation}

	% Herleitung der Einheiten
	\begin{align*}
		[B] = \frac{\text{H/m} \cdot \text{A} \cdot 1}{\text{m}}
		= \frac{\text{H} \cdot \text{A}}{\text{m}^2}
		= \frac{(\text{V} \cdot \text{s / } \colorcancel{\text{A}}{red}) \cdot \colorcancel{\text{A}}{red}}{\text{m}^2}
		= \frac{\text{V} \cdot \text{s}}{\text{m}^2}
		= \text{T (Tesla)}
	\end{align*}
	% Formel für die magnetische Flussdichte
	\vspace{0.5cm}
	\begin{equation}
		B = \frac{\mu \cdot I \cdot N} {l}
	\end{equation}



	% Erklärung der Symbole und Herleitung der Einheiten
	\begin{align*}
		\mu & : \text{Permeabilitätskonstante} \\
		I   & : \text{Elektrische Stromstärke} \\
		N   & : \text{Anzahl der Windungen}    \\
		l   & : \text{Länge des Magnetkreises}
	\end{align*}
}

%%%%%%%%%%%%%%%%%%%%%%%%%%%%%%%%%%%%%%%%%%%%%%%%% Ende - Formeln Induktivität %%%%%%%%%%%%%%%%%%%%%%%%%%%%%%%%%%%%%%%%%%%%%%
\begin{frame}
	\ftx{Die Induktivität}
	\b{
		\begin{equation*}
			[L] = 1\, \text{Henry} = 1\,\text{H} = 1\, \frac{\text{Vs}} {\text{A}}
		\end{equation*}
		\vspace{0.5cm}
		\begin{align*}
			L = \frac{\mu \cdot N^2 \cdot A}{l}
		\end{align*}
	}
	\speech{04_Induktitaet_und_Spule14}{1}
	{Die Induktivität einer Spule hängt neben dem verwendeten magnetisch permeablen Material, der Querschnittsfläche und der Länge, quadratisch von der Wicklungsanzahl. Eine nähere Herleitung ist in Modul 6 zu finden. 
	}
	\silence{4}
\end{frame}
\s{
	Die Induktivität selbst hängt ebenfalls von der Bauform und den verwendeten Materialien ab.
	\begin{equation*}
		[L] = 1\, \text{Henry} = 1\, \text{H} = 1\, \frac{\text{Vs}} {\text{A}}
	\end{equation*}
	\begin{equation}
		L = \frac{\mu \cdot N^2 \cdot A}{l}
	\end{equation}
}

%%%%%%%%%%%%%%%%%%%%%%%%%%%%%%%%%%%%%%%%%%%%%%%%% Anfang - mag. Flussdichte B	 %%%%%%%%%%%%%%%%%%%%%%%%%%%%%%%%%%%%%%%%%%%%%%
% \subsection{Die magnetische Feldstärke}% $\vec{H}$}

% \s{
% 	\begin{equation}
% 		\vec{B} = \mu \cdot \vec{H}
% 	\end{equation}
% }
%%%%%%%%%%%%%%%%%%%%%%%%%%%%%%%%%%%%%%%%%%%%%%%%% Ende - mag. Flussdichte B	 %%%%%%%%%%%%%%%%%%%%%%%%%%%%%%%%%%%%%%%%%%%%%%

%%%%%%%%%%%%%%%%%%%%%%%%%%%%%%%%%%%%%%%%%%%%%%%%% Anfang - Formeln Induktivität	 %%%%%%%%%%%%%%%%%%%%%%%%%%%%%%%%%%%%%%%%%%%%%%

%%%%%%%%%%%%%%%%%%%%%%%%%%%%%%%%%%%%%%%%%%%%%%%%% Ende - Formeln Induktivität	 %%%%%%%%%%%%%%%%%%%%%%%%%%%%%%%%%%%%%%%%%%%%%%


%%%%%%%%%%%%%%%%%%%%%%%%%%%%%%%%%%%%%%%%%%%%%%%%% Anfang - Schaltverhalten Spule	 %%%%%%%%%%%%%%%%%%%%%%%%%%%%%%%%%%%%%%%%%%%%%%
\s{
\begin{frame}
	%\ftx{Schaltverhalten einer Spule: Aufladen}
	\subsection{Schaltverhalten einer Spule}

		Zum besseren Verständnis der Funktion einer Spule wird nun ein zeitlich veränderliches Signal, hier der Einschalt- und Ausschaltvorgang, betrachtet.
		Das grundlegende Prinzip der Induktivität basiert darauf, dass ein zeitlich veränderlicher Strom ein zeitlich veränderliches Magnetfeld erzeugt. Dieses zeitlich veränderliche Magnetfeld
		induziert wiederum einen Induktionsstrom in die Spule selbst. Aufgrund der Lenz'schen Regel wirkt der Induktionsstrom jedoch entgegen seiner Ursache, was einen sprunghaften Anstieg des Stromes verhindert.
		% Das grundlegende Prinzip der Induktivität basiert darauf, dass ein zeitlich veränderlicher Strom ein zeitlich veränderliches Magnetfeld erzeugt.
		% Dieses Magnetfeld induziert wiederum eine Spannung in der Spule, die der Änderung des Stromes entgegenwirkt. Dieser Vorgang wird als Selbstinduktion bezeichnet.
	

	%%%%% Laden der Spule
	%%%%%%%%%%%%%%%%%%%%%%%%%%%%%%%%%%%%%%%%%%%%%%%%% Anfang - Schaltbild zum Schaltverhalten SCHALTER	 %%%%%%%%%%%%%%%%%%%%%%%%%%%%%%%%%%%%%%%%%%%%%%
	%\b{
		 %\begin{figure}[H]
		% 	\centering
		% 	\begin{subfigure}[h]{0.5\textwidth}
		% 		\centering
		% 		\input{Tikz/src/04_Induktivitaet_Spule/Reihenschaltung_Spule_Aufladen}
		% 		\label{fig:Schaltbild_Spule_Schalter}
		% 	\end{subfigure}
			%%%%%%%%%%%%%%%%%%%%%%%%%%%%%%%%%%%%%%%%%%%%%%%%% Ende - Schaltbild zum Schaltverhalten SCHALTER	 %%%%%%%%%%%%%%%%%%%%%%%%%%%%%%%%%%%%%%%%%%%%%%
			%%%%%%%%%%%%%%%%%%%%%%%%%%%%%%%%%%%%%%%%%%%%%%%%% Anfang - Diagramm zum Schaltverhalten SCHALTER	 %%%%%%%%%%%%%%%%%%%%%%%%%%%%%%%%%%%%%%%%%%%%%%
		% 	\begin{subfigure}[b]{0.9\textwidth}
		% 		\centering
		% 		\input{Tikz/src/04_Induktivitaet_Spule/Spule_Aufladen_Verlauf}
		% 		\label{fig:Spannungsverlauf}
		% 	\end{subfigure}
			% \alttext{Zweiteilige Abbildung zum Einschaltvorgang eines RL-Glieds. Oben ist eine Reihenschaltung aus Gleichspannungsquelle Uq, Schalter, Induktivität L und Widerstand R dargestellt. Der Schalter wird bei t0 gleich 5 Tau geöffnet. Pfeile markieren die Spannung uL von t an der Induktivität, die Spannung uR von t am Widerstand sowie den Strom i von t durch den Zweig. Unten zeigt ein Zeitdiagramm die normierten Verläufe: Uq springt bei t gleich 0 auf 100 Prozent und bleibt bis t0 konstant, danach fällt sie auf null. Der Strom i von t steigt exponentiell von 0 auf 100 Prozent an. Die Induktorspannung uL von t startet bei 100 Prozent und fällt exponentiell gegen 0. Bei t gleich Tau sind etwa 63 Prozent für i und 37 Prozent für uL markiert.}
		% 	\label{fig:Gesamtdarstellung}
		% \end{figure}
% 	}% nur Folien
 \end{frame}
}

%\s{
	%\begin{figure}[H]
		% \centering
		% \begin{subfigure}[h]{0.9\textwidth}
		% 	\centering
		% 	\input{Tikz/src/04_Induktivitaet_Spule/Reihenschaltung_Spule_Aufladen}
		% 	\caption{Darstellung der Reihenschaltung bestehend aus einer Induktivität \(L\), einem Widerstand \(R\) und einem Schalter.}
		% 	\label{fig:Schaltbild_Spule_Schalter}
		% \end{subfigure}
		%%%%%%%%%%%%%%%%%%%%%%%%%%%%%%%%%%%%%%%%%%%%%%%%% Ende - Schaltbild zum Schaltverhalten SCHALTER	 %%%%%%%%%%%%%%%%%%%%%%%%%%%%%%%%%%%%%%%%%%%%%%
		%%%%%%%%%%%%%%%%%%%%%%%%%%%%%%%%%%%%%%%%%%%%%%%%% Anfang - Diagramm zum Schaltverhalten SCHALTER	 %%%%%%%%%%%%%%%%%%%%%%%%%%%%%%%%%%%%%%%%%%%%%%
	% 	\begin{subfigure}[b]{0.9\textwidth}
	% 		\centering
	% 		\input{Tikz/src/04_Induktivitaet_Spule/Spule_Aufladen_Verlauf}
	% 		\caption{Spannungs- und Strom verlauf über der Induktivität bei einer Gleichspannung die zum Zeitpunkt $t = 0$ eingeschaltet
	% 			wird und zum Zeitpunkt $t_\mathrm{0}$ abgeschaltet wird.}
	% 		\label{fig:Spannungsverlauf}
	% 	\end{subfigure}
	% 	\caption{\textbf{Der Spannungs- und Strom verlauf über der Induktivität.} Das Diagramm zeigt die Reaktion der Induktivität
	% 		auf die plötzliche Änderung der Spannung und
	% 		veranschaulicht das typische Lade verhalten.}
			% \alttext{Zweiteilige Abbildung zum Einschaltvorgang eines RL-Glieds. Oben ist eine Reihenschaltung aus Gleichspannungsquelle Uq, Schalter, Induktivität L und Widerstand R dargestellt. Der Schalter wird bei t0 gleich 5 Tau geöffnet. Pfeile markieren die Spannung uL von t an der Induktivität, die Spannung uR von t am Widerstand sowie den Strom i von t durch den Zweig. Unten zeigt ein Zeitdiagramm die normierten Verläufe: Uq springt bei t gleich 0 auf 100 Prozent und bleibt bis t0 konstant, danach fällt sie auf null. Der Strom i von t steigt exponentiell von 0 auf 100 Prozent an. Die Induktorspannung uL von t startet bei 100 Prozent und fällt exponentiell gegen 0. Bei t gleich Tau sind etwa 63 Prozent für i und 37 Prozent für uL markiert.}
	% 	\label{fig:Gesamtdarstellung}
	% \end{figure}
%}% Nur Skript
%%%%%%%%%%%%%%%%%%%%%%%%%%%%%%%%%%%%%%%%%%%%%%%%% Ende - Diagramm zum Schaltverhalten SCHALTER	 %%%%%%%%%%%%%%%%%%%%%%%%%%%%%%%%%%%%%%%%%%%%%%

%%%%%%%%%%%%%%%%%%%%%%%%%%%%%%%%%%%%%%%%%%%%%%%%% Anfang - Schaltbild zum Schaltverhalten RECHTECK	 %%%%%%%%%%%%%%%%%%%%%%%%%%%%%%%%%%%%%%%%%%%%%%
\begin{frame}
	\ftx{Schaltverhalten einer Spule: Aufladen und Entladen}
	%%%%%% Laden- entladen der Spule
	\b{
		\begin{figure}[H]
			\centering
			\begin{subfigure}[h]{0.5\textwidth}
				\centering
				\input{Tikz/src/04_Induktivitaet_Spule/Reihenschaltung_Spule_Entladen}
				\label{fig:Schaltbild_Spule_Rechteck}
			\end{subfigure}
			%%%%%%%%%%%%%%%%%%%%%%%%%%%%%%%%%%%%%%%%%%%%%%%%% Ende - Schaltbild zum Schaltverhalten RECHTECK	 %%%%%%%%%%%%%%%%%%%%%%%%%%%%%%%%%%%%%%%%%%%%%%
			%%%%%%%%%%%%%%%%%%%%%%%%%%%%%%%%%%%%%%%%%%%%%%%%% Anfang - Diagramm zum Schaltverhalten RECHTECK	 %%%%%%%%%%%%%%%%%%%%%%%%%%%%%%%%%%%%%%%%%%%%%%
			\uncover<1->{
			\begin{subfigure}[b]{0.9\textwidth}
				\centering
				\input{Tikz/src/04_Induktivitaet_Spule/Spule_aufladen_entladen_folie}
				\label{fig:Spannungsverlauf}
			\end{subfigure}
			\alttext{Zweiteilige Abbildung zum Lade- und Entladeverhalten eines RL-Glieds bei rechteckförmiger Anregung. Oben ist eine Reihenschaltung aus Induktivität L und Widerstand R an einer rechteckförmigen Spannungsquelle u von t dargestellt; Pfeile markieren die Induktorspannung uL von t, die Widerstandsspannung uR von t und den Strom i von t. Unten zeigt ein Zeitdiagramm drei Verläufe: die Eingangsspannung u von t springt auf 100 Prozent und fällt bei 5 Tau auf 0. Der Strom i von t steigt bis 5 Tau exponentiell gegen 100 Prozent und fällt danach exponentiell wieder gegen 0. Die Induktorspannung uL von t startet bei 100 Prozent und fällt bis 5 Tau gegen 0; nach dem Abschalten wird uL von t negativ und nähert sich wieder 0. Markierungen zeigen bei Tau etwa 63 Prozent für i und 37 Prozent für uL; nach dem Abschalten sind bei 6 Tau Punkte auf i und uL markiert.}
			\label{fig:Gesamtdarstellung}
			}
		\end{figure}
	}
	\speech{04_Induktitaet_und_Spule14}{1}
	{Schauen wir uns nun zuerst das auflade und dann das entladeverhalten einer induktivität an. Hierbei treten besonders beim entladen effekte auf, die nicht direkt intuitiv sind.
	Wir haben als Grundlage eine Spannungsquelle welche zum Zeitpunkt Null eingeschaltet wird, und nach 5 Tau, was genau Tau ist, kommt am ende, wieder auf null Volt geschaltet wird. 
	Wichtig: Eine Spannungquelle die Null Volt hat ist ein Kurzschluss. Das ist vorerst verwirrend, aber wenn man drüber nachdenkt, die Quelle erzwingt, unabhängig vom Strom eine Spannung von  null Volt, was gleichzusetzen mit einem idealen Leiter, also einem kurzschluss ist. Denn es fällt keine Spannung bei einem beliebigen Strom ab.
	}
	\silence{3}
	\speech{04_Induktitaet_und_Spule14}{2}
	{Schalten wir nun die Spannung an, und fokussieren uns auch nur auf den bereich bis 5 Tau.
	Wir einnern uns, durch eine induktivität kann sich der Strom nicht Sprunghaft ändern, was bedeutet, die induktivität zeigt lässt zu beginn keinen Strom fließen. Demzufolge fällt keine Spannung am Widerstand R ab, da kein Strom fließt. Die gesamte Spannung fällt also über die induktivität ab.
	Der Strom durch eine induktivität ist ja das integral der Spannung über die Zeit, heißt, die induktivität leitet mit der zeit immer mehr Strom, was auch an der roten kurve im Graphen zu sehen ist.
	Durch den Stromfluss fällt auch über den Widerstand R eine Spannung ab, die Spannung an der induktivität verringert sich um exakt diesen Wert. Das geht so lange, bis die gesamte Quellspannung U über den Widerstand abfällt. Über die Induktivität fällt dann keine Spannung mehr ab, die integriert werden kann, der Strom verändert sich nicht weiter.
	}
	\silence{3}
	\speech{04_Induktitaet_und_Spule14}{3}
	{Nun wird die Quelle auf Null gesetzt.
	Das bedeutet, der Strom wird nun nicht mehr durch die Quelle getrieben. Man könnte meinen der Strom fließt dann auch nicht mehr weiter, aber das ist bei der induktivität ja nicht möglich. Die Energie im Feld der induktivität treibt den Strom weiter durch die Schaltung, bis die Energie vom Widerstand in thermische Energie umgewandelt ist.
	}
	\silence{3}
	\speech{04_Induktitaet_und_Spule14}{4}
	{Mit der Spannung an der induktivität passiert aber etwas eigenartiges. Sie dreht sich um. Was zuerst unintuitiv wirkt, aber sinn macht, wenn wir die Quelle als kurzschluss sehen. denn damit der Strom in gleicher Richtung fließt muss die urpsrüngliche Quellspannung am Widerstand R anliegen. Und damit ist die Spannung anderherum als am Anfang.
	Hierbei ist ganz wichtig, das die Spannung durch die Induktivität induziert wird. Wenn wir statt eines konstanten Widerstandes nun einen Schalter hätten, der plötzlich einen unendlich hohen Widerstand hat, würde die Induzierte Spannung so hoch werden, bis der urpsrüngliche Strom wieder fließt. Dabei können auch lichtbögen im Schalter entstehen.
	}
	\silence{3}
	\speech{04_Induktitaet_und_Spule14}{5}
	{Die Anfangs erwähnte Zeitkonstante Tau ist ähnlich bei bei der Kapazität die zeit, dies es braucht bis sich der Strom oder auch die Spannung um 63 Prozent zum anfangswert verändert hat. Dabei ist es egal, ob die Spannung Zu oder abnimmt.}
\end{frame}
\s{
	\begin{figure}[H]
		\centering
		\begin{subfigure}[h]{0.85\textwidth}
			\centering
			\input{Tikz/src/04_Induktivitaet_Spule/Reihenschaltung_Spule_Entladen}
			\caption{Darstellung der Reihenschaltung bestehend aus einer Induktivität \(L\) und einem Widerstand \(R\) an einer rechteckförmigen Spannungsquelle.}
			\label{fig:Schaltbild_Spule_Rechteck}
		\end{subfigure}
		%%%%%%%%%%%%%%%%%%%%%%%%%%%%%%%%%%%%%%%%%%%%%%%%% Ende - Schaltbild zum Schaltverhalten RECHTECK	 %%%%%%%%%%%%%%%%%%%%%%%%%%%%%%%%%%%%%%%%%%%%%%
		%%%%%%%%%%%%%%%%%%%%%%%%%%%%%%%%%%%%%%%%%%%%%%%%% Anfang - Diagramm zum Schaltverhalten RECHTECK	 %%%%%%%%%%%%%%%%%%%%%%%%%%%%%%%%%%%%%%%%%%%%%%
		\begin{subfigure}[b]{0.9\textwidth}
			\centering
			\input{Tikz/src/04_Induktivitaet_Spule/Spule_Entladen_Verlauf}
			\caption{Spannungs- und Strom verlauf über der Induktivität bei einer rechteckförmigen Spannung.}
			\label{fig:Spannungsverlauf}
		\end{subfigure}
		\caption{\textbf{Spannungs- und Stromverlauf über der Induktivität.} Das Diagramm zeigt die Reaktion der Induktivität auf eine plötzliche Änderung
			der Spannung und veranschaulicht das typische Lade- und Entladeverhalten der Induktivität.}
			\alttext{Zweiteilige Abbildung zum Lade- und Entladeverhalten eines RL-Glieds bei rechteckförmiger Anregung. Oben ist eine Reihenschaltung aus Induktivität L und Widerstand R an einer rechteckförmigen Spannungsquelle u von t dargestellt; Pfeile markieren die Induktorspannung uL von t, die Widerstandsspannung uR von t und den Strom i von t. Unten zeigt ein Zeitdiagramm drei Verläufe: die Eingangsspannung u von t springt auf 100 Prozent und fällt bei 5 Tau auf 0. Der Strom i von t steigt bis 5 Tau exponentiell gegen 100 Prozent und fällt danach exponentiell wieder gegen 0. Die Induktorspannung uL von t startet bei 100 Prozent und fällt bis 5 Tau gegen 0; nach dem Abschalten wird uL von t negativ und nähert sich wieder 0. Markierungen zeigen bei Tau etwa 63 Prozent für i und 37 Prozent für uL; nach dem Abschalten sind bei 6 Tau Punkte auf i und uL markiert.}
		\label{fig:Gesamtdarstellung}
	\end{figure}
}
%%%%%%%%%%%%%%%%%%%%%%%%%%%%%%%%%%%%%%%%%%%%%%%%% Ende - Diagramm zum Schaltverhalten RECHTECK	 %%%%%%%%%%%%%%%%%%%%%%%%%%%%%%%%%%%%%%%%%%%%%%
%%%%%%%%%%%%%%%%%%%%%%%%%%%%%%%%%%%%%%%%%%%%%%%%% Ende - Schaltverhalten Spule	 %%%%%%%%%%%%%%%%%%%%%%%%%%%%%%%%%%%%%%%%%%%%%%
\s{
	In der praktischen Anwendung ermöglicht die Induktivität die Konstruktion von Komponenten wie Transformatoren, Drosseln und Induktivitäten,
	die in einer Vielzahl von elektrischen und elektronischen Geräten zu finden sind. Von der Filterung von Signalen über die Energieübertragung bis
	hin zur Steuerung elektromagnetischer Interferenzen ist die Induktivität ein Schlüsselelement in der modernen Elektrotechnik.
}
%%%%%%%%%%%%%%%%%%%%%%%%%%%%%%%% Anfang - Beispiel Rechnung Induktivität %%%%%%%%%%%%%%%%%%%%%%%%%%%%%%%
% \begin{frame}
% 	\ftx{Beispielrechnung Induktivität}

% 	\begin{bsp}{Berechnungs der Induktivität}{}
% 		\begin{itemize}
% 			\item Berechnung der Induktivität $L$
% 			\item Berechnung der magnetischen Flussdichte $\vec{B}$
% 			\item Gespeicherte Energie im Magnetfeld $E_\mathrm{m}$ ?
% 		\end{itemize}
% 	\end{bsp}

% \end{frame}
%%%%%%%%%%%%%%%%%%%%%%%%%%%%%%%% Ende - Beispiel Rechnung Induktivität %%%%%%%%%%%%%%%%%%%%%%%%%%%%%%%
%%%%%%%%%%%%%%%%%%%%%%%%%%%%%%%%%%%%%%%%%%%%%%%%% Ende - Induktivität %%%%%%%%%%%%%%%%%%%%%%%%%%%%%%%%%%%%%%%%%%%%%%
